\subsection{Biological setting}
AMPs are produced across all domains of life. The canonical family --- short,
net-cationic, amphipathic peptides such as magainins, LL-37, and the
defensins --- partitions into microbial membranes and kills by
permeabilization, but the known universe also includes anionic AMPs,
glycine-rich peptides, proline-rich peptides that enter cells and inhibit
translation, and larger disulfide-stabilized folds. This mechanistic variety
is the central recognition challenge: there is no single sequence signature,
and any single-feature rule (charge, hydrophobic moment, length) fails on
whole families. It is also why we treat sensitivity numbers honestly: a model
trained on a corpus dominated by short cationic peptides should not be
expected to score disulfide-locked anionic families highly, but per-family sensitivity has not yet been measured in this study.

\subsection{Contributions}
\begin{enumerate}
\item An auditable, exact-10-mer cluster-split benchmark with a subsequently
identified source-review confound:
24{,}628 UniProt KW-0929 positives (mostly unreviewed) against 24{,}628 length-matched
reviewed negatives, clustered by shared exact 10-mers into 12{,}110
components and split 80/10/10 by cluster, with all fetch URLs, hashes, and
timestamps recorded.
\item A paired comparison of two neural architectures and two composition
baselines under this split, with bootstrap intervals and permutation tests
computed by one shared harness: PepCNN attains descriptive AUROC $0.921$ [0.913, 0.929] on the confounded
held-out sample, ahead of the tested baselines; the graph model is a reported negative
result.
\item An APD6 2024 positive-versus-reviewed-UniProt-negative comparison after exact-10-mer filtering:
PepCNN scores AUROC $0.926$ [0.918, 0.934] and sensitivity $0.601$ at a
5\% false-positive rate, but review-status confounding in training limits
its biological interpretation.
\item A simulated-annealing designer with a fixed-temperature detailed-balance
proof producing ranked sequence candidates, and a discovery
screen over 852 reviewed uncharacterized proteins producing 25 ranked
hypotheses, all clearly labeled as in-silico.
\item A reproducible software artifact: the \texttt{pepdesign} package, multiple study scripts, a 27-test hermetic suite, and reproducible
benchmark and candidate tables derived from checked-in result files.
\end{enumerate}

\subsection{Scope}
This paper is bounded to an AMP benchmark provenance audit and its exploratory recognition, design and screening work. No activity against specific strains, toxicity, structural mechanism or wet-lab outcome is inferred from binary labels and unassayed sequences.
